% Chapter Template

\chapter{Ensayos y resultados} % Main chapter title

\label{Chapter4} % Change X to a consecutive number; for referencing this chapter elsewhere, use \ref{ChapterX}
En este capítulo se exponen los datos con los que se trabajó y se expande sobre los resultados obtenidos en base a los experimentos comentados en el capítulo anterior. Se los exponen en el formato en que se investigaron, de más simple a más complejo. Se finaliza con el objetivo principal de este trabajo que son los modelos de transformers.

\section{Datos utilizados}

Como ya se mencionó antes, en este trabajo se utilizaron precios de cierre de Bitcoin en relación a USDT entre 2020-01-01 y 2025-11-07. En la tabla \ref{tab:buy_and_hold} se ve el rendimiento de bitcoin durante el periodo analizado.

\begin{table}[h]
  \centering
  \caption[Resumen Buy and Hold]{Resumen de rendimiento de la estrategia Buy and Hold entre 2020-01-01 y 2025-11-07.}
  \begin{tabular}{l c}
    \toprule
    \textbf{Métrica} & \textbf{Valor} \\
    \midrule
    Retorno Total & 31995,25\% \\
    Volatilidad anualizada & 70,72\% \\
    Ratio de Sharpe & 1,16 \\
    Drawdown máximo & -83,27\% \\
    \bottomrule
  \end{tabular}
  \label{tab:buy_and_hold}
\end{table}

También se ilustra como las barras de volumen afectan a la serie de tiempo original (figura \ref{fig:bars-vs-1m}). Se aprecia como la forma general sigue presente pero, en cierta manera, un poco más suavizada. Se pueden ver que los movimientos pequeños se condensan en un solo punto.

\begin{figure}[H]
  \centering
  \includegraphics[width=1.0\textwidth]{"Figures/bars-vs-1m.png"}
  \caption{Comparativa de barras de volumen y precio de cierre con frecuencia de 1m.}
  \label{fig:bars-vs-1m}
\end{figure}


La figura \ref{fig:bars-volume-1m} ilustra, además, cómo el agregado por barras de volumen cambia la velocidad de muestreo con respecto un incremento en la varianza. Esto se ejemplifica bien si se compara el periodo que va del  9 de enero y al 13 de enero, con el que va del 13 al 21 del mismo mes.
Tienen la misma cantidad de puntos (seis) pero el segundo período cubre el doble de tiempo.



\begin{figure}[H]
  \centering
  \includegraphics[width=1.0\textwidth]{"Figures/bars-volume-1m.png"}
  \caption{Superposición de precio de cierre cada 1m, barras de volumen a 50 000 unidades de BTC, y volumen de BTC.}
  \label{fig:bars-volume-1m}
\end{figure}


\section{Pruebas con algoritmos de trading}
\label{sec:pruebasTrading}
En esta sección se hizo un análisis sobre las estrategias basadas en algoritmos de trading. Esto es, se codificaron reglas fijas basadas en distintos indicadores y, en base a ellas, se determinó una señal de compra o venta para la ocasión actual.

Configuración: CV purgado de 5 pliegues con 1\% de embargo, 0.1\% de comisión, optimización Optuna en datos de 1 minuto de BTCUSDT.

Como se observa en la tabla \ref{tab:optimizacion_f1}, el \texttt{BollingerBandsClassifier} tiene el techo de rendimiento más alto. Por otro lado, la tabla \ref{tab:sharpe_dist} muestra que el \texttt{RSIDivergenceClassifier} es la estrategia más confiable: presenta el Sharpe medio más alto y una varianza extremadamente baja.



\vspace{1em}

% --- TABLA 1 ---
\begin{table}[h]
    \centering
    \caption[Resultados de Optimización]{Resultados de la mejor optimización basados en el promedio pesado de F1 en validación cruzada ($avg\_cv\_weighted\_f1$).}
    \begin{tabular}{l l c l}
        \toprule
        \textbf{Rank} & \textbf{Clasificador} & \textbf{Mejor F1} & \textbf{Mejores Parámetros} \\
        \midrule
        1 & BollingerBandsClassifier & 0.5108 & bb\_window=21, bb\_std=1.77 \\
        2 & RSIDivergenceClassifier  & 0.5079 & rsi\_window=9, divergence\_period=23 \\
        3 & MaCrossoverClassifier    & 0.5012 & short=6, long=7 \\
        4 & MultiIndicatorClassifier & 0.4986 & fast\_sma=42, slow\_sma=227, ... \\
        5 & MACDClassifier           & 0.4985 & fast=5, slow=6, signal=17 \\
        6 & SmaCrossClassifier       & 0.4981 & n1=47, n2=199 \\
        \bottomrule
    \end{tabular}
    \label{tab:optimizacion_f1}
\end{table}

% --- TABLA 2 ---
\begin{table}[h]
    \centering
    \caption[Distribución de Sharpe]{Distribución del ratio de Sharpe a nivel de prueba (184 pruebas totales).}
    \begin{tabular}{l c c c c c c}
        \toprule
        \textbf{Clasificador} & \textbf{N} & \textbf{Media} & \textbf{Std} & \textbf{Min} & \textbf{Median} & \textbf{Max} \\
        \midrule
        RSIDivergenceClassifier  & 32 & 0.5072 & 0.0008 & 0.5051 & 0.5077 & 0.5079 \\
        BollingerBandsClassifier & 30 & 0.5043 & 0.0221 & 0.3882 & 0.5093 & 0.5108 \\
        SmaCrossClassifier       & 30 & 0.4963 & 0.0019 & 0.4909 & 0.4968 & 0.4981 \\
        MACDClassifier           & 41 & 0.4957 & 0.0019 & 0.4924 & 0.4957 & 0.4985 \\
        MaCrossoverClassifier    & 30 & 0.4956 & 0.0024 & 0.4919 & 0.4950 & 0.5012 \\
        MultiIndicatorClassifier & 21 & 0.4899 & 0.0270 & 0.3731 & 0.4973 & 0.4986 \\
        \bottomrule
    \end{tabular}
    \label{tab:sharpe_dist}
\end{table}

La tabla \ref{tab:consistencia} cuantifica esta consistencia y muestra el porcentaje de pruebas que supera distintos umbrales de Sharpe.

% --- TABLA 3 ---
\begin{table}[h]
    \centering
    \caption[Consistencia Sharpe]{Consistencia: porcentaje de pruebas que superan los umbrales de Sharpe establecidos.}
    \begin{tabular}{l c c c c c}
        \toprule
        \textbf{Clasificador} & \textbf{$>$0.49} & \textbf{$>$0.495} & \textbf{$>$0.500} & \textbf{$>$0.505} & \textbf{$>$0.510} \\
        \midrule
        RSIDivergenceClassifier  & 100\% & 100\% & 100\% & 100\% & 0\% \\
        BollingerBandsClassifier & 97\%  & 97\%  & 97\%  & 77\%  & 37\% \\
        SmaCrossClassifier       & 100\% & 77\%  & 0\%   & 0\%   & 0\% \\
        MACDClassifier           & 100\% & 71\%  & 0\%   & 0\%   & 0\% \\
        MaCrossoverClassifier    & 100\% & 50\%  & 7\%   & 0\%   & 0\% \\
        MultiIndicatorClassifier & 81\%  & 71\%  & 0\%   & 0\%   & 0\% \\
        \bottomrule
    \end{tabular}
    \label{tab:consistencia}
\end{table}

La tabla \ref{tab:sensibilidad} resume la sensibilidad de los parámetros clave para cada clasificador.

% --- TABLA 4 ---
\begin{table}[h]
    \centering
    \caption[Sensibilidad de Parámetros]{Sensibilidad de parámetros (correlación con $avg\_cv\_sharpe$).}
    \begin{tabular}{l l c l}
        \toprule
        \textbf{Clasificador} & \textbf{Parámetro} & \textbf{Corr} & \textbf{Interpretación} \\
        \midrule
        SmaCrossClassifier       & n1 & +0.89 & Ventana lenta más larga $\rightarrow$ mejor \\
        MaCrossoverClassifier    & long\_window & -0.58 & Ventanas cortas funcionan mejor \\
        BollingerBandsClassifier & bb\_std & -0.40 & Se prefieren bandas estrechas \\
        MACDClassifier           & slow\_span & -0.92 & Spans cortos fuertemente preferidos \\
        RSIDivergenceClassifier  & rsi\_window & -0.76 & Ventana RSI corta es mejor (best: 9) \\
        \bottomrule
    \end{tabular}
    \label{tab:sensibilidad}
\end{table}


Conclusión: RSIDivergence y BollingerBands son los dos mejores exponentes. El primero es la opción segura para implementación por su baja varianza (ver tabla \ref{tab:consistencia}), mientras que las bandas de Bollinger ofrecen mayor recompensa a cambio de un riesgo superior.

\section{Pruebas con algoritmos estadísticos clásicos}
En esta sección se hizo un análisis sobre las estrategias basadas en algoritmos tradicionales de series de tiempo.

Configuración: CV purgado de 3 pliegues, 1\% de embargo, datos de 1h de BTCUSDT desde 2020-01-01. Se evaluaron 5 familias de modelos.

Como se detalla en la tabla \ref{tab:mejores_modelos}, el \texttt{ARIMAClassifier} es el mejor modelo estadístico evaluado: alcanza un F1 máximo de 0,5004. Sin embargo, su alta varianza ($\text{std} = 0,118$) y sus modos de falla catastróficos lo hacen poco confiable en la práctica, como se evidencia en la tabla \ref{tab:distribucion_f1}.
\begin{table}[H]
  \centering
  \caption[Resultados de Optimización]{Mejor resultado de optimización por familia de modelo basado en el promedio de F1 en validación cruzada ($avg\_cv\_f1$).}
  \begin{tabular}{l l c c l}
    \toprule
    \textbf{Rank} & \textbf{Clasificador} & \textbf{Mejor F1} & \textbf{n\_trials}\\
    \midrule
    1 & ARIMAClassifier        & 0.5004 & 30  \\
    2 & ARIMAXGARCHClassifier  & 0.4954 & 10  \\
    3 & SARIMAClassifier       & $\sim$0.4862 & 4 \\
    4 & KalmanARIMAClassifier  & 0.4184 & 10  \\
    \bottomrule
  \end{tabular}
  \label{tab:mejores_modelos}
\end{table}

Los modelos de series temporales muestran una inestabilidad estructural. La tabla \ref{tab:sensibilidad_estadistica} revela que el parámetro threshold es el factor determinante: las familias ARIMA solo funcionan cuando se configuran para predecir retornos esperados extremadamente pequeños (umbral $\sim 1e-6$).

Además de esta limitación, los modelos ARIMA presentan modos de falla catastróficos. Las altas desviaciones estándar reportadas en la tabla \ref{tab:distribucion_f1} (0,118 para ARIMA, 0,116 para GARCH y 0,165 para Kalman) son un claro indicador de  inestabilidad. Las combinaciones deficientes de parámetros producen puntuaciones F1 cercanas a cero, lo que los hace poco confiables para aplicaciones prácticas. Como evidencia de esto, ARIMA solo logra un F1 $\geq 0,48$ en el 47\% de las pruebas, mientras que Kalman nunca logra alcanzar esa marca.

\begin{table}[H]
  \centering
  \caption[Distribución de F1]{Distribución de la métrica F1 a través de todas las pruebas realizadas por clasificador.}
  \begin{tabular}{l c c c c c c}
    \toprule
    \textbf{Clasificador} & \textbf{Media} & \textbf{Std} & \textbf{Min} & \textbf{Median} & \textbf{Max} & \textbf{$\ge$0.50} \\
    \midrule
    SARIMAClassifier       & 0.4605 & 0.0457 & 0.3922 & 0.4819 & 0.4862 & 0\% \\
    ARIMAClassifier        & 0.4382 & 0.1183 & 0.0214 & 0.4768 & 0.5004 & 3\% \\
    ARIMAXGARCHClassifier  & 0.3855 & 0.1159 & 0.1880 & 0.4062 & 0.4954 & 0\% \\
    KalmanARIMAClassifier  & 0.2476 & 0.1653 & 0.0094 & 0.3117 & 0.4184 & 0\% \\
    \bottomrule
  \end{tabular}
  \label{tab:distribucion_f1}
\end{table}

Queda en evidencia que, en este contexto, añadir complejidad perjudica el rendimiento. Al contrastar los resultados de la tabla \ref{tab:mejores_modelos}, un modelo ARIMA (6 parámetros) supera a los más complejos ARIMAX-GARCH (8 parámetros) y Kalman-ARIMA (6 parámetros). Cada nueva capa de complejidad reduce el F1 promedio y aumenta la varianza de los resultados. El filtro de Kalman, en particular, parece degradar la señal original del modelo ARIMA en lugar de potenciarla.

\begin{table}[H]
  \centering
  \caption[Sensibilidad de Parámetros]{Análisis de sensibilidad: Correlación entre hiperparámetros clave y el rendimiento (F1).}
  \begin{tabular}{l l c l}
    \toprule
    \textbf{Modelo} & \textbf{Parámetro} & \textbf{Corr} & \textbf{Interpretación} \\
    \midrule
    ARIMA-family        & threshold       & -0.93 & El umbral cercano a cero es crítico. \\
    ARIMAXGARCH         & threshold       & -0.99 & Umbrales altos degradan el modelo. \\
    ARIMAClassifier     & lookback\_len   & -0.37 & Se prefiere historial corto (720h). \\
    ARIMAXGARCH         & $q$ (ARIMA MA)  & +0.85 & Orden MA alto es significativamente mejor. \\
    \bottomrule
  \end{tabular}
  \label{tab:sensibilidad_estadistica}
\end{table}

En cuanto a la parametrización, la mejor configuración de ARIMA converge de manera muy ajustada. Las cinco mejores pruebas de este modelo comparten exactamente los mismos valores de base ($p=4$, $d=0$, $q=1$, $\text{refit}=720\text{h}$, $\text{lookback}=720\text{h}$), y difieren únicamente en el umbral. Esto sugiere que el orden autorregresivo y de media móvil de ARIMA está bien identificado para este activo, pero que el espacio de búsqueda inicial para el umbral fue demasiado amplio. 

\vspace{1em}

Conclusión: podemos afirmar con seguridad que se logró un resultado subóptimo aunque esperable. Los algoritmos estadísticos no supieron diferenciar la escasa señal del preponderante ruido presente en los mercados financieros. En resultados concretos el mejor f1 no pudo superar el de un intento aleatorio y en sharpe ratio quedó muy por debajo de la base del buy and hold.

En el detalle de los modelos estudiados, todos los algoritmos estadísticos de series temporales (ARIMA, GARCH y Kalman) resultan estructuralmente inestables para este conjunto de datos. Exigen una calibración del umbral extremadamente precisa e, incluso en el mejor de los escenarios (\texttt{ARIMAClassifier} con $p=4$, $d=0$, $q=1$), apenas logran superar el rendimiento de un clasificador aleatorio.

\section{Pruebas con algoritmos de machine learning}

Aquí es se realizaron las primeras tareas significativas de ingeniería de características y de etiquetado  ya que, previamente, las estrategias recibían solamente el precio de cierre. También en esta etapa es donde se definió nuestra estrategia de baseline, la inspirada en el trabajo de Palazzo.

Se evaluaron dos configuraciones de pipeline sobre klines de BTCUSDT en velas de 1 minuto (enero 2020 -- octubre 2025, 3.033.173 filas), resumidas en la tabla \ref{tab:ml_experiments}. El etiquetado es común a todos los pipelines: barrera triple binaria sobre barras de volumen, con la condición $\text{bar\_return}[t+1] \geq \text{bar\_return}[t] + \text{intra\_bar\_std}[t] \times \tau$, agregadas a partir de ticks de 1 minuto. La validación cruzada empleó Purged K-Fold con embargo (\texttt{n\_splits=3}, \texttt{pct\_embargo=0.01}); la métrica principal es el F1 macro promedio entre pliegues (\texttt{avg\_cv\_f1}).

\begin{table}[H]
    \centering
    \caption[Experimentos de optimización ML]{Experimentos de optimización de machine learning.}
    \begin{tabular}{l l l l}
        \toprule
        \textbf{ID} & \textbf{Experimento} & \textbf{Modelo} & \textbf{Datos} \\
        \midrule
        59 & \texttt{Palazzo\_XGBoost\_Optimization}    & XGBoost (CUDA) & 1-min consolidado \\
        61 & \texttt{AutoGluon\_Palazzo\_Optimization}  & AutoGluon      & 1-min consolidado \\
        \bottomrule
    \end{tabular}
    \label{tab:ml_experiments}
\end{table}

\subsection{Distribución de F1 por algoritmo}

La tabla \ref{tab:ml_f1_dist} resume la distribución del F1 macro promedio para cada algoritmo a lo largo de todos los trials de optimización. AutoGluon Palazzo alcanza un rendimiento similar al de XGBoost Palazzo con menos trials pero un costo computacional considerablemente mayor.

\begin{table}[H]
    \centering
    \caption[Distribución F1 ML]{Distribución del F1 macro promedio por algoritmo (optimización con Optuna).}
    \begin{tabular}{l c c c c c c c}
        \toprule
        \textbf{Algoritmo} & \textbf{$n$} & \textbf{mín} & \textbf{p25} & \textbf{mediana} & \textbf{p75} & \textbf{máx} & \textbf{std} \\
        \midrule
        XGBoost Palazzo  & 60 & 0.584 & 0.606 & 0.612 & 0.623 & \textbf{0.647} & 0.015 \\
        AutoGluon Palazzo & 20 & 0.561 & 0.606 & 0.616 & 0.618 & \textbf{0.639} & 0.016 \\
        \bottomrule
    \end{tabular}
    \label{tab:ml_f1_dist}
\end{table}

\subsection{Hallazgos principales}

\begin{itemize}
\item AutoGluon iguala a XGBoost con un costo computacional diez veces mayor: AutoGluon Palazzo (20 trials, mejor F1=0.639) se ubica a 0.008 del mejor XGBoost Palazzo (0.647). El mejor modelo seleccionado por AutoGluon fue \texttt{LightGBMXT\_BAG\_L1}. Dado que cada trial de AutoGluon insume aproximadamente diez veces más tiempo, XGBoost resulta la opción práctica para esta escala de datos.

\item La reducción de dimensionalidad por PCA no aporta beneficio: ningún experimento se benefició de PCA; todas las optimizaciones convergieron a \texttt{use\_pca=False}. La compresión PCA elimina señal discriminativa en este conjunto de características.

\item \texttt{tau} y \texttt{volume\_threshold} convergen consistentemente: como se observa en la tabla \ref{tab:ml_label_params}, todos los pipelines convergen a $\tau > 0.98$, lo que corresponde a barreras amplias y etiquetas menos frecuentes pero de mayor calidad. El umbral de volumen varía considerablemente entre experimentos (28.522--62.855), lo que sugiere que la calidad del etiquetado es sensible a la granularidad de las barras.

\begin{table}[H]
    \centering
    \caption[Parámetros de etiquetado óptimos]{Parámetros de etiquetado óptimos por pipeline.}
    \begin{tabular}{l c c}
        \toprule
        \textbf{Pipeline} & \textbf{\texttt{tau} óptimo} & \textbf{\texttt{volume\_threshold} óptimo} \\
        \midrule
        XGBoost Palazzo   & 1.261 & 28.522 \\
        AutoGluon Palazzo & 0.980 & 32.034 \\
        \bottomrule
    \end{tabular}
    \label{tab:ml_label_params}
\end{table}

\item Hiperparámetros de XGBoost Palazzo: la tabla \ref{tab:ml_xgb_params} muestra los principales hiperparámetros del mejor trial. El modelo converge a árboles poco profundos (\texttt{max\_depth=4}) y a una tasa de aprendizaje muy baja (0.00328), coherente con un conjunto de datos de más de 3 millones de filas que requiere alta regularización.

\end{itemize}

\begin{table}[H]
    \centering
    \caption[Hiperparámetros óptimos XGBoost]{Hiperparámetros óptimos de XGBoost Palazzo.}
    \begin{tabular}{l c}
        \toprule
        \textbf{Parámetro} & \textbf{XGBoost Palazzo} \\
        \midrule
        \texttt{max\_depth}     & 4       \\
        \texttt{learning\_rate} & 0.00328 \\
        \texttt{n\_estimators}  & 373     \\
        \bottomrule
    \end{tabular}
    \label{tab:ml_xgb_params}
\end{table}

\section{Pruebas con Chronos}

En esta sección se describen las pruebas realizadas con nuestro modelo transformer de elección, Chronos. Como ya se dijo se eligió este modelo por ser la destilación estado del arte actual en series de tiempo y por su fácil acceso a través de librerías como AutoGluon.

Se exploraron tres familias arquitectónicas de pipelines basadas en Chronos, enumeradas en la tabla \ref{tab:chronos_experiments}. Todos los experimentos utilizan los mismos datos que la etapa anterior (enero 2020 -- octubre 2025, 3.033.173 filas) con barras de volumen fijas en \texttt{volume\_threshold=50.000} BTC y $\tau=0.7$. La validación cruzada empleó Purged K-Fold con embargo (\texttt{n\_splits=3}, \texttt{pct\_embargo=0.01}) para las familias de embeddings, y \texttt{pct\_embargo=0.05} para predicción directa.

\begin{table}[H]
    \centering
    \caption[Experimentos Chronos]{Experimentos con pipelines Chronos organizados por familia arquitectónica.}
    \begin{tabular}{l l l c}
        \toprule
        \textbf{Exp} & \textbf{Pipeline} & \textbf{Familia} & \textbf{Trials} \\
        \midrule
        73 & \texttt{ChronosFeaturePipeline}             & Embedding + clasificador & 10 \\
        71 & \texttt{ChronosBoltFeaturePipeline}         & Embedding + clasificador & 10 \\
        65 & \texttt{Chronos2FeaturePipeline}            & Embedding + clasificador & 10 \\
        67 & \texttt{FinetunedChronosFeaturePipeline}    & Embedding + clasificador & 10 \\
        69 & \texttt{PalazzoChronosBinaryClassificationPipeline} & Predicción directa & 5 \\
        22 & \texttt{ChronosMetaLabelingPipeline}        & Meta-etiquetado          & 8  \\
        \bottomrule
    \end{tabular}
    \label{tab:chronos_experiments}
\end{table}

\subsection{Familia 1: embeddings como características}

Cada pipeline de esta familia extrae representaciones densas de una ventana de precios de cierre recientes mediante un encoder de Chronos, luego concatena esos embeddings con las características tabulares de \texttt{PalazzoXGBoostPipeline} y clasifica la matriz combinada con AutoGluon. La métrica principal es \texttt{avg\_cv\_f1} (F1 macro promedio entre pliegues).

La tabla \ref{tab:chronos_f1_dist} resume la distribución de F1 para cada pipeline a lo largo de los trials de optimización. \texttt{ChronosFeaturePipeline} (T5) presenta la mediana y el máximo más altos, aunque con mayor varianza. \texttt{ChronosBoltFeaturePipeline} muestra la varianza más baja de todas (std=0.008), lo que refleja mayor estabilidad en los embeddings de BOLT.

\begin{table}[H]
    \centering
    \caption[Distribución F1 Chronos embeddings]{Distribución del F1 macro promedio por pipeline de embeddings (optimización con Optuna).}
    \begin{tabular}{l c c c c c c c}
        \toprule
        \textbf{Pipeline} & \textbf{$n$} & \textbf{mín} & \textbf{p25} & \textbf{mediana} & \textbf{p75} & \textbf{máx} & \textbf{std} \\
        \midrule
        \texttt{ChronosFeaturePipeline} (T5)     & 10 & 0.573 & 0.579 & 0.605 & 0.624 & \textbf{0.633} & 0.022 \\
        \texttt{ChronosBoltFeaturePipeline}       & 10 & 0.598 & 0.607 & 0.609 & 0.613 & \textbf{0.628} & 0.008 \\
        \texttt{Chronos2FeaturePipeline}          &  3 & 0.604 & 0.610 & 0.615 & 0.618 & \textbf{0.621} & 0.007 \\
        \texttt{FinetunedChronosFeaturePipeline}  &  4 & 0.598 & 0.599 & 0.604 & 0.610 & \textbf{0.611} & 0.005 \\
        \bottomrule
    \end{tabular}
    \label{tab:chronos_f1_dist}
\end{table}

\subsection{Familia 2: predicción directa}

\texttt{PalazzoChronosBinaryClassificationPipeline} (exp 69) afina Chronos-2-small sobre etiquetas $\{+1, -1\}$ y convierte la salida continua en señal binaria por umbral en cero. La optimización con Optuna (5 trials, métrica \texttt{avg\_cv\_score}) arrojó los resultados de la tabla \ref{tab:chronos_direct}.

\begin{table}[H]
    \centering
    \caption[Resultados predicción directa Chronos]{Resultados de optimización de \texttt{PalazzoChronosBinaryClassificationPipeline} (exp 69).}
    \begin{tabular}{c c c c c}
        \toprule
        \textbf{Trial} & \textbf{avg\_cv\_score} & \textbf{test\_macro\_f1} & \textbf{tau} & \textbf{volume\_threshold} \\
        \midrule
        1 (mejor) & \textbf{0.818} & \textbf{0.819} & 0.910 & 45.988 \\
        3         & 0.816          & 0.819          & 1.002 & 50.565 \\
        0         & 0.813          & 0.814          & 0.871 & 50.817 \\
        2         & 0.807          & 0.809          & 1.218 & 45.718 \\
        4         & 0.803          & 0.804          & 1.134 & 53.145 \\
        \bottomrule
    \end{tabular}
    \label{tab:chronos_direct}
\end{table}

\subsection{Familia 3: meta-etiquetado}

\texttt{ChronosMetaLabelingPipeline} (exp 22) entrena un modelo primario con embeddings T5-tiny más características tabulares y luego un meta-modelo sobre si el primario acertará. En 8 corridas con \texttt{amazon/chronos-t5-tiny} el modelo primario obtuvo \texttt{baseline\_f1\_weighted=0.655} de forma consistente. Sin embargo, el meta-filtro produjo valores de \texttt{metalabeling\_f1\_weighted} entre 0.571 y 0.657: en 6 de las 8 corridas degradó el F1 del modelo primario.

\subsection{Hallazgos principales}

\begin{itemize}

\item El fine-tuning directo con Chronos-2 es el mejor resultado obtenido: \texttt{PalazzoChronosBinaryClassificationPipeline} alcanza \texttt{avg\_cv\_score=0.818}, significativamente por encima de todos los pipelines de embeddings (0.611--0.633) y del meta-etiquetado (0.655 en F1 ponderado).

\item T5 tiene el mayor potencial dentro de los embeddings: \texttt{ChronosFeaturePipeline} obtiene el F1 más alto (0.633 en exp 73) y también la mayor varianza (std=0.022), lo que indica mayor sensibilidad a los hiperparámetros. El modelo ganador en la optimización fue \texttt{CatBoost\_BAG\_L1\_FULL} con \texttt{chronos-t5-small}, ventana de 96 y preset \texttt{high\_quality}.

\item BOLT es estable pero con techo inferior: \texttt{ChronosBoltFeaturePipeline} muestra std=0.008, la varianza más baja de todos los pipelines evaluados. Su máximo (0.628) queda por debajo de T5. Dentro de la familia BOLT, bolt-small supera a bolt-mini con presupuestos de tiempo cortos.

\item Chronos-2 y el fine-tuning por pliegue no superan a T5 ni a BOLT: \texttt{Chronos2FeaturePipeline} alcanza F1=0.621 y \texttt{FinetunedChronosFeaturePipeline} apenas 0.611, pese al costo computacional significativo del ajuste fino por pliegue. La señal binaria sobre barras de volumen resulta demasiado ruidosa para mejorar las representaciones pre-entrenadas en 300 segundos.

\item El meta-filtro no mejora al modelo primario: agregar un meta-modelo de AutoGluon sobre las predicciones del clasificador primario no produce una mejora consistente. La arquitectura de meta-etiquetado impone un compromiso precisión-recall que perjudica el F1 macro en el conjunto de prueba más frecuentemente de lo que lo beneficia.

\end{itemize}

%% \subsection{Mejores configuraciones encontradas}

%% A continuación se detallan los hiperparámetros completos del mejor trial de cada pipeline. La familia de embeddings emplea \texttt{avg\_cv\_f1} como métrica; la predicción directa emplea \texttt{avg\_cv\_score}: ambas no son directamente comparables.

%% \textbf{ChronosFeaturePipeline (T5) -- F1: 0.633} (exp 73, trial 8)

%% \begin{verbatim}
%% chronos_model_name=amazon/chronos-t5-small   window_size=96   stride=1
%% tau=0.7 (fijo)   volume_threshold=50000   use_pca=False
%% AutoGluon: presets=high_quality  time_limit=600s
%% best_model=CatBoost_BAG_L1_FULL
%% \end{verbatim}

%% \textbf{ChronosBoltFeaturePipeline -- F1: 0.628} (exp 71, trial 9)

%% \begin{verbatim}
%% chronos_model_name=amazon/chronos-bolt-mini   window_size=32   stride=1
%% tau=0.7 (fijo)   volume_threshold=50000   use_pca=False
%% AutoGluon: presets=medium_quality  time_limit=600s
%% \end{verbatim}

%% \textbf{Chronos2FeaturePipeline -- F1: 0.621} (exp 65, trial 0)

%% \begin{verbatim}
%% chronos_model_name=autogluon/chronos-2-small   window_size=96   stride=1
%% tau=0.7 (fijo)   volume_threshold=50000   use_pca=False
%% AutoGluon: presets=best_quality  time_limit=600s
%% \end{verbatim}

%% \textbf{FinetunedChronosFeaturePipeline -- F1: 0.611} (exp 67, trial 7)

%% \begin{verbatim}
%% chronos_model_name=amazon/chronos-t5-tiny   window_size=160   stride=1
%% tau=0.7 (fijo)   volume_threshold=50000
%% AutoGluon: presets=medium_quality  time_limit=600s
%% \end{verbatim}

%% \textbf{PalazzoChronosBinaryClassificationPipeline -- avg\_cv\_score: 0.818} (exp 69, trial 1)

%% \begin{verbatim}
%% chronos_model_name=autogluon/chronos-2-small
%% tau=0.910   volume_threshold=45988
%% \end{verbatim}


\section{Selección del mejor modelo}
